\chapter{衍射}
\problem{4-1}{在菲涅耳圆孔衍射实验中，圆孔半径 $2.0\,\mathrm{mm}$，光源离圆孔 $2.0\,\mathrm m$，波长 $0.5\,\mu\mathrm m$，当接收屏幕由很远的地方向圆孔靠近时，求：（1）前三次出现中心亮斑（强度极大）的位置；（2）前三次出现中心暗斑（强度极小）的位置。}
\answerarea
\problem{4-2}{在菲涅耳圆孔衍射实验中，光源距离圆孔 $1.5\,\mathrm m$，波长 $0.63\,\mu\mathrm m$，接收屏幕与圆孔距离 $6.0\,\mathrm m$，圆孔半径从 $0.5\,\mathrm{mm}$ 开始逐渐扩大。求：（1）最先的两次出现中心亮斑时圆孔的半径；（2）最先的两次出现中心暗斑时圆孔的半径。}
\answerarea
\clearpage
\problem{4-3}{用直刀口将点光源的波面遮住一半（直边衍射），几何阴影边缘点上的光强比自由传播时小多少倍？}
\answerarea
\problem{4-4}{求圆孔中露出 $1.5$ 个半波带时衍射场中心强度与自由传播时强度之比。}
\answerarea
\clearpage
\problem{4-5}{用平行光照明衍射屏，如对波面作如本题图所示的下几种方式的遮挡，求轴上点的光强与自由传播时之比（图中标出的是该处到场点的光程，其中 $b$ 是中心到场点的光程）。\fig[0.85]{4_5}}
\answerarea
\problem{4-6}{若一个菲涅耳波带片只将前五个偶数半波带遮挡，其余地方都开放，求衍射场中心强度与自由传播时之比。}
\answerarea
\clearpage
\problem{4-7}{若一个菲涅耳波带片将前50个奇数半波带遮挡，其余地方都开放，求衍射场中心强度与自由传播时之比。}
\answerarea
\problem{4-8}{菲涅耳波带片第一个半波带的半径 $\rho_1=5.0\,\mathrm{mm}$。（1）用波长 $\lambda=1.06\,\mu\mathrm m$ 的单色平行光照明，求主焦距；（2）若要求主焦距为 $25\,\mathrm{cm}$，需将此波带片缩小多少倍？}
\answerarea
\clearpage
\problem{4-9}{如何制作一张满足以下要求的波带片：（1）它在 $400.0\,\mathrm{nm}$ 紫光照明下的主焦距为 $80\,\mathrm{cm}$；（2）主焦点光强是自由传播时的 $10^4$ 倍左右。}
\answerarea
\problem{4-10}{一菲涅耳波带片对 $900.0\,\mathrm{nm}$ 的红外光主焦距为 $30\,\mathrm{cm}$，改用 $632.8\,\mathrm{nm}$ 的氦氖激光照明，主焦距为多少？}
\answerarea
\clearpage
\problem{4-11}{如本题图，平行光以 $\theta_0$ 角斜入射在宽度为 $a$ 的单缝上。试证明：
\begin{enumerate}\item 夫琅禾费衍射的强度公式基本不变（忽略倾斜因子），即
\[I_\theta=I_0\left(\frac{\sin\alpha}{\alpha}\right)^2,\qquad \alpha=\frac{\pi a}{\lambda}(\sin\theta-\sin\theta_0),\]
式中 $I_0$ 为零级中心强度，只不过 $\alpha$ 的定义与正入射不同；\item 零级中心的位置在几何光学像点处；\item 零级斑半角宽度为 $\Delta\theta=\lambda/(a\cos\theta_0)$。\end{enumerate}\fig[0.9]{4_11}
\editorial{对照PDF，第（3）问分母同样印为 $a\sin\theta_0$；按第（1）问在 $\theta=\theta_0$ 附近展开，分母应为 $a\cos\theta_0$，此处据此校正。}}
\answerarea
\clearpage
\problem{4-12}{估算下列情形光束在介质界面上反射和折射时反射光束和折射光束的衍射发散角。设界面的线度为 $1\,\mathrm{cm}$，光波波长为 $0.6\,\mu\mathrm m$，折射率为 $1.5$：（1）平行光正入射；（2）入射角为 $75^\circ$；（3）入射角为 $89^\circ$（掠入射）。}
\answerarea
\problem{4-13}{试用巴比涅原理证明：互补的衍射屏产生的夫琅禾费衍射图样相同。}
\answerarea
\clearpage
\problem{4-14}{衍射细丝测径仪就是把单缝夫琅禾费衍射装置中的单缝用细丝代替。今测得零级衍射斑的宽度（两个一级暗纹间的距离）为 $1\,\mathrm{cm}$，求细丝的直径。已知光波波长 $0.63\,\mu\mathrm m$，透镜焦距 $50\,\mathrm{cm}$。}
\answerarea
\problem{4-15}{一对双星的角间隔为 $0.05''$。（1）需要多大口径的望远镜才能分辨它们？（2）此望远镜的角放大率应设计为多少才比较合理？}
\answerarea
\clearpage
\problem{4-16}{一台天文望远镜的口径为 $2.16\,\mathrm m$，由这一数据你能进一步获得关于它在光学性能方面的哪些知识？}
\answerarea
\problem{4-17}{一台显微镜，已知其 $\mathrm{N.A.}=1.32$，物镜焦距 $f'_o=1.91\,\mathrm{mm}$，目镜焦距 $f'_E=50\,\mathrm{mm}$。求：（1）最小分辨距离；（2）有效放大率；（3）光学筒长。}
\answerarea
\clearpage
\problem{4-18}{用一架照相机在离地面 $200\,\mathrm{km}$ 的高空拍摄地面上的物体，如果要求它能分辨地面上相距 $1\,\mathrm m$ 的两点，照相机的镜头至少要多大？设镜头的几何像差已很好地消除，感光波长为 $400.0\,\mathrm{nm}$。}
\answerarea
\problem{4-19}{已知地月距离约为 $3.8\times10^5\,\mathrm{km}$，用口径为 $1\,\mathrm m$ 的天文望远镜能分辨月球表面两点的最小距离是多少？}
\answerarea
\clearpage
\problem{4-20}{已知日地距离约为 $1.5\times10^8\,\mathrm{km}$，要求分辨太阳表面相距 $20\,\mathrm{km}$ 的两点，望远镜的口径至少需有多大？}
\answerarea
\problem{4-21}{用坐标纸绘制 $N=2$、$d=3a$ 的夫琅禾费衍射强度分布曲线，横坐标取 $\sin\theta$，至少画到第7级主极大，并计算第0级和第1级主极大与单缝主极大之比。}
\answerarea
\clearpage
\problem{4-22}{用坐标纸绘制 $N=6$、$d=1.5a$ 的夫琅禾费衍射强度分布曲线，横坐标取 $\sin\theta$，至少画到第4级主极大，并计算第4级主极大与单缝主极大之比。}
\answerarea
\problemfig[0.37]{4-23}{导出斜入射时夫琅禾费多缝衍射强度分布公式：
\[I_\theta=I_0\left(\frac{\sin\alpha'}{\alpha'}\right)^2\left(\frac{\sin N\beta'}{\sin\beta'}\right)^2,\]
其中 $\alpha'=\dfrac{\pi a}{\lambda}(\sin\theta-\sin\theta_0)$，$\beta'=\dfrac{\pi d}{\lambda}(\sin\theta-\sin\theta_0)$。$\theta_0$ 为入射线与光轴的夹角（见本题图）。}{4_23}
\answerarea
\clearpage
\problem{4-24}{导出斜入射时主极大位置公式、第 $k$ 级主极大的半角宽度公式及缺级情况，并注意与正入射情况作比较。}
\answerarea
\problem{4-25}{有三条平行狭缝，宽度都是 $a$，缝距分别为 $d$ 和 $2d$（见本题图）。证明正入射时夫琅禾费衍射强度分布公式为
\[I_\theta=I_0\left(\frac{\sin\alpha}{\alpha}\right)^2[3+2(\cos2\beta+\cos4\beta+\cos6\beta)],\]
其中 $\alpha=\pi a\sin\theta/\lambda$，$\beta=\pi d\sin\theta/\lambda$。\fig{4_25}}
\answerarea
\clearpage
\problem{4-26}{导出不等宽双缝的夫琅禾费衍射强度分布公式，缝宽分别为 $a$ 和 $2a$，缝距 $d=3a$（见本题图）。\fig{4_26}}
\answerarea
\problemfig[0.43]{4-27}{有 $2N$ 条平行狭缝，缝宽相同都是 $a$，缝间不透明部分的宽度作周期性变化 $a$、$3a$、$a$、$3a$、$\ldots$（见本题图）。求下列各种情形中的衍射强度分布：（1）遮住偶数缝；（2）遮住奇数缝；（3）全开放。}{4_27}
\answerarea
\clearpage
\problem{4-28}{波长为 $650.0\,\mathrm{nm}$ 的红光谱线，被观察发现它是双线。如果在 $9\times10^4$ 条刻线光栅的第3级光谱中恰好能分辨此双线，求其波长差。}
\answerarea
\problem{4-29}{若要50条/mm的光栅在第2级光谱中能分辨钠双线 $\lambda_1=589.0\,\mathrm{nm}$ 和 $\lambda_2=589.6\,\mathrm{nm}$，光栅宽度应选多少？}
\answerarea
\clearpage
\problem{4-30}{绿光 $500.0\,\mathrm{nm}$ 正入射在光栅常数为 $2.5\times10^{-4}\,\mathrm{cm}$、宽度为 $3\,\mathrm{cm}$ 的光栅上，聚光镜的焦距为 $50\,\mathrm{cm}$。（1）求第1级光谱的线色散；（2）求第1级光谱中能分辨的最小波长差；（3）设光栅最多能看到第几级光谱？}
\answerarea
\problem{4-31}{一束白光正入射在600条/mm的光栅上，第1级可见光谱末端与第2级光谱前端之间的角间隔有多少？}
\answerarea
\clearpage
\problem{4-32}{国产31W1型1米平面光栅摄谱仪的技术数据表如下，试根据这些数据来计算：
\begin{center}\begin{tabular}{l|l}\hline
物镜焦距&$1050\,\mathrm{mm}$\\光栅刻划面积&$60\,\mathrm{mm}\times40\,\mathrm{mm}$\\
闪耀波长&$635.0\,\mathrm{nm}$（1级）\\刻线&1200条/mm\\色散&$0.8\,\mathrm{nm/mm}$\\理论分辨率&72000（1级）\\\hline\end{tabular}\end{center}
（1）该摄谱仪能分辨的谱线间隔的最小值为多少？（2）该摄谱仪的角色散本领为多少〔以 $\mathrm{nm}/('')$ 为单位〕？（3）光栅的闪耀角为多大？闪耀方向与光栅平面的法线方向成多大角度？}
\answerarea
\problem{4-33}{底边长度为 $6\,\mathrm{cm}$ 的棱镜，在光波长为 $0.6\,\mu\mathrm m$ 附近能分辨的最小波长间隔为多少？设制棱镜材料的色散 $\mathrm dn/\mathrm d\lambda$ 值为 $0.4\times10^{-4}\,\mathrm{nm}^{-1}$，来估算。}
\answerarea
\clearpage
\problem{4-34}{根据以下数据比较光栅、棱镜、法--珀干涉仪三者的分光本领：（1）分辨本领；（2）色散本领；（3）自由光谱范围。
\par 光栅宽度 $D=5\,\mathrm{cm}$，刻线密度 $1/d=600$条/mm；棱镜底边 $b=5\,\mathrm{cm}$，顶角 $\alpha=60^\circ$，折射率 $n=1.51$，色散率 $\mathrm dn/\mathrm d\lambda=0.6\times10^{-4}/\mathrm{nm}$；法--珀腔长 $h=5\,\mathrm{cm}$，反射率 $R=0.99$。}
\answerarea
\clearpage
